\section{Absolute Zero：让模型自己提出下一道题}

\subsection{训练任务本身成为数据瓶颈}

\videofigure{figures/task-data-bottleneck.jpg}{当前 reasoning training 依赖人类专家编写 prompt、答案或 reasoning trace，能力越强越难找到足够前沿任务。}{00:23:46--00:24:38}

传统 RLVR 的 verifier 可以自动评分，但问题集合仍由人静态提供。模型进步后，简单题失去梯度，过难题又全错；训练需要持续处于 current frontier 的任务。

\videofigure{figures/proposer-solver.jpg}{Absolute Zero 用同一个模型同时扮演 proposer 与 solver，在可执行代码环境中自生成任务并自我求解。}{00:24:38--00:25:24}

\videofigure{figures/absolute-zero-loop.jpg}{Proposer 根据 learnability reward 生成任务，solver 根据 accuracy reward 解题，两种信号联合更新模型。}{00:25:24--00:26:18}

\begin{importantbox}{Absolute Zero 消除的是外部 task data，不是环境}
系统仍需要代码执行器提供可验证反馈。所谓 zero data 指不依赖人类预先编写的问题集，而不是在没有环境和 reward 的真空中学习。
\end{importantbox}

\subsection{Deduction、Abduction 与 Induction}

\videofigure{figures/proposer-modes.jpg}{Proposer 通过 deduction、abduction 与 induction 三种代码任务模式扩大自生成课程。}{00:26:18--00:27:10}

三类任务可理解为：

\begin{itemize}
    \item \textbf{Deduction}：给程序与输入，预测执行输出；
    \item \textbf{Abduction}：给程序与输出，反推出满足条件的输入；
    \item \textbf{Induction}：给若干输入输出与自然语言描述，推断程序或规则。
\end{itemize}

这些任务都能由执行器判定，因此适合自动生成、求解和验证；三种方向又要求不同 reasoning，有助于避免课程过于单一。

\subsection{Learnability reward：不要太简单，也不要完全不可能}

设同一任务上 solver 的经验成功率为 $\bar p$。课程描述的简化 proposer reward 为：

$$
r_{\mathrm{learn}}=
\begin{cases}
0, & \bar p=0,\\
1-\bar p, & \bar p>0.
\end{cases}
$$

\begin{itemize}
    \item $\bar p=1$：任务已掌握，learnability 低；
    \item $\bar p=0$：任务可能无效或远超当前能力，也不给奖励；
    \item $0<\bar p<1$：任务处在能力边界，既有成功样本又有改进空间。
\end{itemize}

\begin{knowledgebox}{这是自动化 curriculum learning}
随着 solver 变强，过去困难任务变简单，proposer 必须持续增加复杂度。课程不再由固定题库定义，而由模型当前失败边界动态生成。
\end{knowledgebox}

\subsection{Task validation 与 buffer management}

\videofigure{figures/task-validation.jpg}{任务进入训练前要通过可执行性、安全性与确定性检查，并写入带成功统计的 task buffer。}{00:28:10--00:29:28}

验证包括程序能否运行、是否触发危险操作、相同输入是否确定地产生相同输出。Buffer 保存程序、输入、输出、成功率和参考任务，使 proposer 能从历史课程取样并生成不同但可学习的新任务。

\begin{warningbox}{Task proposer 也会 reward hack}
它可能生成利用解释器漏洞的题、伪装成困难但实际无意义的题，或重复稍作改写的旧任务。需要 sandbox、去重、复杂度约束和独立 validation。
\end{warningbox}

\subsection{结果与开放问题}

\videofigure{figures/absolute-zero-takeaways.jpg}{Absolute Zero 在无人工 prompt 数据下获得强代码 reasoning，并呈现任务复杂度与多样性随训练增长的 emergent curriculum。}{00:29:28--00:33:34}

课程强调 proposer 与 solver 带有轻微对抗性：proposer 想找到 solver 还不会但能学会的问题，solver 则不断压缩这一区域。结果表明自生成任务可以超过用大量专家例子训练的基线。

\begin{importantbox}{真正的开放性来自任务分布也能变化}
只在固定题库上重复 RL，最多逼近该分布的上限；让模型改变问题本身，才有机会持续扩展学习边界。但任务有效性与现实相关性随之成为新 verifier 问题。
\end{importantbox}

\subsection{本章小结}

Absolute Zero 把 synthetic data 从“自己生成答案”推进到“自己生成题目与答案”。代码执行环境提供即时验证，使 proposer--solver 能形成动态课程；向开放现实领域推广时，环境真实性和 reward 成为核心限制。

